\subsection*{CU-31: Responder una solicitud de amistad}
\addcontentsline{toc}{subsection}{CU-31: Responder una solicitud de amistad}
\label{cu-31}

\begin{fichacu}{CU-31}{Responder una solicitud de amistad}
  \crudFila{Responsable}{Ángel Gabriel Aguilar Hernández}
  \crudFila{Actor principal}{Jugador destinatario}
  \crudFila{Actores interesados}{Jugador solicitante}
  \crudFila{Actores de sistema}{Servidor de partidas, Base de datos}
  \crudFila{Prototipos}{5c, 2b}
  \crudFila{Objetivo}{Permitir al Jugador aceptar o rechazar las solicitudes de amistad que recibió, y cancelar las que envió y siguen sin respuesta.}
  \crudFila{Disparador}{El Jugador abre la pestaña de solicitudes en la \texttt{GUI\_\allowbreak{}Friends}, o selecciona la notificación de una solicitud nueva.}
  \textbf{Precondiciones} & \cuCelda{\begin{cureglas}
      \item \textbf{\mbox{PRE-01}} El Jugador tiene una \textit{Sesion} abierta y su token es válido.
      \item \textbf{\mbox{PRE-02}} El \textit{Usuario} tiene \texttt{tipo\_\allowbreak{}cuenta} en \texttt{REGISTRADA}.
      \item \textbf{\mbox{PRE-03}} El Servidor de partidas está en ejecución y tiene conexión disponible con la Base de datos.
    \end{cureglas}} \\ \hline
  \textbf{Flujo normal} & \cuCelda{\begin{cuenum}[start=1]
      \item El Jugador abre la pestaña de solicitudes. \textit{(Ver \mbox{FA-06})}
      \item El SJM envía el mensaje \texttt{LIST\_\allowbreak{}FRIEND\_\allowbreak{}REQUESTS} con el token de sesión. \textit{(Ver \mbox{EX-02}, \mbox{EX-03}, \mbox{EX-04})}
      \item El Servidor de partidas recupera las \textit{Solicitud} en las que el \textit{Usuario} es destinatario y las que envió, descartando las caducadas (\mbox{RN-05}). \textit{(Ver \mbox{EX-01})}
      \item El Servidor de partidas responde con las solicitudes entrantes y enviadas, con el \texttt{nickname}, el icono y el elo del modo clásico de cada jugador.
      \item El SJM despliega las dos listas: entrantes, con ``Aceptar'' y ``Rechazar'', y enviadas, con ``Cancelar''. \textit{(Ver \mbox{FA-01})}
      \item El Jugador selecciona ``Aceptar'' en una solicitud entrante. \textit{(Ver \mbox{FA-02}, \mbox{FA-03})}
      \item El SJM envía el mensaje \texttt{FRIEND\_\allowbreak{}REQUEST\_\allowbreak{}ACCEPT} con el identificador de la solicitud.
      \item El Servidor de partidas verifica que la \textit{Solicitud} exista y que el \textit{Usuario} sea su destinatario (\mbox{RN-01}). \textit{(Ver \mbox{FA-04})}
    \end{cuenum}} \\
   & \cuCelda{\begin{cuenum}[start=9]
      \item El Servidor de partidas verifica que ninguno de los dos jugadores haya alcanzado el límite de cien amigos (\mbox{RN-02}). \textit{(Ver \mbox{FA-05})}
      \item El Servidor de partidas verifica que no haya surgido un \textit{Bloqueo} entre ambos desde que se envió. \textit{(Ver \mbox{FA-04})}
      \item El Servidor de partidas inicia una transacción, crea la \textit{Amistad} con el identificador menor en \texttt{id\_\allowbreak{}usuario\_\allowbreak{}a}, el mayor en \texttt{id\_\allowbreak{}usuario\_\allowbreak{}b} y la \texttt{fecha\_\allowbreak{}amistad}, y elimina la \textit{Solicitud} (\mbox{RN-03}). \textit{(Ver \mbox{EX-01})}
      \item El Servidor de partidas confirma la transacción y notifica al solicitante, si tiene sesión abierta, con \texttt{FRIEND\_\allowbreak{}REQUEST\_\allowbreak{}ACCEPTED}.
      \item El Servidor de partidas responde con \texttt{FRIEND\_\allowbreak{}REQUEST\_\allowbreak{}ACCEPT\_\allowbreak{}OK}.
      \item El SJM retira la solicitud de la lista y añade al jugador a la lista de amigos.
      \item Fin del caso de uso.
    \end{cuenum}} \\ \hline
  \textbf{Flujos alternos} & \cuCelda{\textbf{\mbox{FA-01}: No hay solicitudes}\par
    \begin{cuenum}[start=1]
      \item El SJM muestra un estado vacío en cada pestaña sin solicitudes, con la acción ``añadir amigo'', que continúa en el \mbox{CU-30}.
      \item Fin del caso de uso.
    \end{cuenum}} \\
   & \cuCelda{\textbf{\mbox{FA-02}: El Jugador rechaza la solicitud}\par
    \begin{cuenum}[start=1]
      \item El Jugador selecciona ``Rechazar'' y el SJM envía \texttt{FRIEND\_\allowbreak{}REQUEST\_\allowbreak{}DECLINE}.
      \item El Servidor de partidas verifica que el \textit{Usuario} sea el destinatario y elimina la \textit{Solicitud}.
      \item El Servidor de partidas \textbf{no} notifica al solicitante (\mbox{RN-04}).
      \item El SJM retira la solicitud de la lista.
      \item Fin del caso de uso.
    \end{cuenum}} \\
   & \cuCelda{\textbf{\mbox{FA-03}: El Jugador cancela una solicitud que envió}\par
    \begin{cuenum}[start=1]
      \item El Jugador selecciona ``Cancelar'' en una solicitud enviada y el SJM envía \texttt{FRIEND\_\allowbreak{}REQUEST\_\allowbreak{}CANCEL}.
      \item El Servidor de partidas verifica que el \textit{Usuario} sea el solicitante y elimina la \textit{Solicitud}.
      \item El SJM retira la solicitud de la lista de enviadas; si el destinatario la tenía abierta, desaparece de la suya en su siguiente recarga.
      \item Fin del caso de uso.
    \end{cuenum}} \\
   & \cuCelda{\textbf{\mbox{FA-04}: La solicitud ya no existe o hay un bloqueo}\par
    \begin{cuenum}[start=1]
      \item El Servidor de partidas no encuentra la \textit{Solicitud}, porque el solicitante la canceló, caducó o se respondió desde otro dispositivo, o encuentra un \textit{Bloqueo} entre ambos.
      \item Si hay bloqueo, el Servidor de partidas elimina la \textit{Solicitud} sin crear la \textit{Amistad}.
      \item El Servidor de partidas responde con \texttt{FRIEND\_\allowbreak{}REQUEST\_\allowbreak{}NOT\_\allowbreak{}FOUND}, igual en todos los casos.
      \item El SJM retira la solicitud de la lista con un aviso breve de que ya no está disponible.
      \item Fin del caso de uso.
    \end{cuenum}} \\
   & \cuCelda{\textbf{\mbox{FA-05}: Uno de los dos alcanzó el límite de amigos}\par
    \begin{cuenum}[start=1]
      \item El Servidor de partidas responde con \texttt{FRIEND\_\allowbreak{}REQUEST\_\allowbreak{}FAILED} y el motivo \texttt{LIMITE\_\allowbreak{}ALCANZADO}, indicando si es el propio Jugador o el solicitante.
      \item El Servidor de partidas conserva la \textit{Solicitud}, para que pueda aceptarse si se libera espacio antes de que caduque.
      \item El SJM muestra la \texttt{GUI\_\allowbreak{}MessageWarning} con el motivo; si el límite es propio, propone revisar la lista de amigos.
      \item Fin del caso de uso.
    \end{cuenum}} \\
   & \cuCelda{\textbf{\mbox{FA-06}: El Jugador responde desde la notificación}\par
    \begin{cuenum}[start=1]
      \item El Jugador selecciona ``Aceptar'' o ``Rechazar'' directamente en el aviso emergente de solicitud nueva.
      \item El flujo continúa en el paso 7 del FN o en el paso 1 del \mbox{FA-02}, sin abrir la pestaña.
    \end{cuenum}} \\ \hline
  \textbf{Excepciones} & \cuCelda{\textbf{\mbox{EX-01}: Fallo al consultar o escribir en la base de datos}\par
    \begin{cuenum}[start=1]
      \item El Servidor de partidas revierte la transacción, de modo que no quede una \textit{Amistad} creada con la \textit{Solicitud} todavía pendiente, ni la solicitud eliminada sin la amistad.
      \item El Servidor de partidas registra la excepción con un identificador de incidencia y responde con \texttt{SERVER\_\allowbreak{}ERROR}.
      \item El SJM muestra la \texttt{GUI\_\allowbreak{}MessageError} con el identificador y conserva la solicitud en la lista.
      \item Fin del caso de uso.
    \end{cuenum}} \\
   & \cuCelda{\textbf{\mbox{EX-02}: Se agota el tiempo de espera de la respuesta}\par
    \begin{cuenum}[start=1]
      \item El SJM no reenvía la respuesta y recarga las listas, que dirán si se aplicó.
      \item Fin del caso de uso.
    \end{cuenum}} \\
   & \cuCelda{\textbf{\mbox{EX-03}: La conexión se interrumpe durante el intercambio}\par
    \begin{cuenum}[start=1]
      \item El SJM muestra la barra de conexión perdida y recarga las listas al recuperarla.
      \item Fin del caso de uso.
    \end{cuenum}} \\
   & \cuCelda{\textbf{\mbox{EX-04}: El mensaje recibido está malformado}\par
    \begin{cuenum}[start=1]
      \item El SJM descarta el mensaje, cierra la conexión y registra en la bitácora local el contenido truncado.
      \item El SJM muestra la \texttt{GUI\_\allowbreak{}MessageError} indicando que la respuesta no pudo interpretarse.
      \item Fin del caso de uso.
    \end{cuenum}} \\ \hline
  \textbf{Postcondiciones} & \cuCelda{\begin{cureglas}
      \item \textbf{\mbox{POST-01}} Si el caso termina por el FN, existe una \textit{Amistad} entre ambos jugadores y ya no existe la \textit{Solicitud}.
      \item \textbf{\mbox{POST-02}} Si el caso termina por el \mbox{FA-02} o el \mbox{FA-03}, la \textit{Solicitud} no existe y no se creó ninguna \textit{Amistad}.
      \item \textbf{\mbox{POST-03}} No existen dos \textit{Amistad} para el mismo par de jugadores, en ningún orden.
      \item \textbf{\mbox{POST-04}} Ninguno de los dos jugadores supera los cien amigos.
    \end{cureglas}} \\ \hline
  \textbf{Reglas de negocio} & \cuCelda{\begin{cureglas}
      \item \textbf{\mbox{RN-01}} Solo el destinatario puede aceptar o rechazar una solicitud, y solo el solicitante puede cancelarla.
      \item \textbf{\mbox{RN-02}} El límite de cien amigos se vuelve a comprobar al aceptar, en los dos jugadores, porque pudo alcanzarse después de enviada la solicitud.
      \item \textbf{\mbox{RN-03}} Una \textit{Amistad} se guarda una sola vez por par, con el identificador menor en \texttt{id\_\allowbreak{}usuario\_\allowbreak{}a}. La tabla declara esa condición con una restricción, de modo que A–B y B–A no puedan coexistir aunque falle la aplicación.
      \item \textbf{\mbox{RN-04}} Rechazar no se notifica al solicitante: para él, la solicitud simplemente deja de estar pendiente.
      \item \textbf{\mbox{RN-05}} Una solicitud sin responder caduca a los treinta días.
      \item \textbf{\mbox{RN-06}} Responder una solicitud elimina la fila de \textit{Solicitud} en todos los casos: la solicitud describe un estado presente, no un hecho que haya que conservar.
    \end{cureglas}} \\ \hline
  \textbf{Observación} & \cuCelda{\textit{Sobre la restricción del par ordenado.}\par
    Guardar la amistad dos veces, una en cada sentido, haría más sencillas algunas consultas, pero duplicaría cada amistad y abriría la puerta a que exista en un sentido y no en el otro. Guardarla una sola vez con el par ordenado exige, a cambio, que las consultas busquen al jugador en cualquiera de las dos columnas. Es un costo pequeño y fijo a cambio de que la base de datos, y no la aplicación, garantice que la amistad es recíproca.} \\ \hline
  \textbf{Datos que toca} & \cuCelda{\begin{cudatos}
      \item \textit{Sesion}: lee por token \textit{(paso 2)}
      \item \textit{Solicitud}: lee las entrantes y enviadas \textit{(pasos 3, 8)}
      \item \textit{Solicitud}: elimina \textit{(pasos 11, \mbox{FA-02}, \mbox{FA-03}, \mbox{FA-04})}
      \item \textit{Usuario}: lee \texttt{nickname} e \texttt{id\_\allowbreak{}icono} de cada jugador \textit{(paso 4)}
      \item \textit{EstadisticaModo}: lee el elo del modo clásico \textit{(paso 4)}
      \item \textit{Amistad}: lee, para contar el límite \textit{(paso 9)}
      \item \textit{Amistad}: crea \textit{(paso 11)}
      \item \textit{Bloqueo}: lee \textit{(paso 10)}
    \end{cudatos}} \\ \hline
\end{fichacu}
\clearpage
